% EDA %
\newpage
\section{Exploratory Data Analysis}
\thispagestyle{plain}
\vspace{0.5em}

% ══════════════════════════════════════════════════════════════
% Section 1: Introduction & Dataset Overview
% ══════════════════════════════════════════════════════════════
\subsection{Dataset Overview}

\subsubsection{Dataset Description}

The training dataset (\texttt{train\_data.csv}) contains \textbf{375,734 rows} and \textbf{246 columns} (including the target variable), covering the period from September 2014 to August 2016. Each row represents one location (with anonymized latitude and longitude) and one start date of a 14-day forecast window.

\subsubsection{Task Definition}

The prediction task is a supervised regression problem: given a set of meteorological, climatic, and geographic features observed at or before the start date, predict the 14-day mean 2-meter temperature at each location. The target variable, \texttt{contest-tmp2m-14d\_\_tmp2m}, is defined as $T_{\text{target}} = (T_{\max} + T_{\min}) / 2$ over the 14-day forecast window. The test set follows the same structure as the training set, but the true target values are withheld for evaluation. Model performance is assessed using root mean squared error (RMSE) between the predicted and observed temperatures.

\subsection{Feature Categorization}

\begin{table}[H]
\centering
\caption{High-level categorization of features in the dataset.}
\label{tab:feature_categories}
\begin{tabular}{@{}lcp{7.5cm}@{}}
\toprule
\textbf{Category} & \textbf{No. Features} & \textbf{Brief Description} \\
\midrule
NMME Temperature Forecasts     & 40  & Latest monthly forecasts and weighted averages (weeks 3--4 and 5--6) of 2\,m temperature from individual NMME models plus ensemble means \\
NMME Precipitation Forecasts   & 40  & Weighted average precipitation rate forecasts for weeks 3--4 and 5--6 from NMME models \\
Contest Meteorological Variables & 15  & 14-day aggregates/averages of SLP, geopotential height, winds, humidity, precipitable water, precipitation, etc. \\
SST \& Sea Ice (PCA)           & 20  & PCA components of sea surface temperature and sea ice concentration from NOAA OISST \\
Climate Indices                & 5   & MJO (phase and amplitude), MEI (ENSO index, rank, phase) \\
2010 Baseline Wind Field PCA   & 120 & PCA components of various wind and geopotential height fields from 2010 baseline \\
Geographic \& Static Features  & 4   & Latitude, longitude, elevation, K\"{o}ppen--Geiger climate classification \\
Others                         & 1   & Start date (for deriving seasonal and temporal patterns) \\
\bottomrule
\end{tabular}
\end{table}


\newpage

% ══════════════════════════════════════════════════════════════
% Section 2: Data Quality Overview
% ══════════════════════════════════════════════════════════════
\section{Data Quality Overview}

\subsection{Target Variable Distribution}

% Figure 1 — Left panel: histogram with mean/median lines; Right panel: top 20 missing columns

\begin{figure}[H]
    \centering
    \includegraphics[width=\textwidth]{figures/fig1_data_quality.png}
    \caption{Data Quality Overview. Left: Distribution of the target variable \texttt{contest-tmp2m-14d} with mean and median markers. Right: Top 20 columns ranked by missing value percentage.}
    \label{fig:data_quality}
\end{figure}

The target variable exhibits an approximately symmetric and near-normal distribution with mean $\approx 11.9\,^\circ\text{C}$ and median $\approx 12.3\,^\circ\text{C}$. The slightly lower mean relative to the median indicates a mild left skew. Most values cluster between $-5\,^\circ\text{C}$ and $30\,^\circ\text{C}$, suggesting a typical temperature range for the U.S.\ forecast period.

\subsection{Missing Value Analysis}

The right panel lists the top columns ranked by missing value percentage. Only 8 columns exhibit missingness, with the highest rates around 4--5\%. These columns belong predominantly to the \textbf{CCSM30} and \textbf{CCSM3} model families (e.g., \texttt{nmme0-tmp2m-34w\_\_ccsm30}, \texttt{nmme-tmp2m-56w\_\_ccsm3}), indicating a systematic issue with those particular NMME model outputs. The overall data quality is high, with the vast majority of features being fully complete.



\newpage

% ══════════════════════════════════════════════════════════════
% Section 3: Temporal Analysis
% ══════════════════════════════════════════════════════════════
\section{Temporal Analysis}

\subsection{Monthly Temperature Distribution}

\begin{figure}[H]
    \centering
    \includegraphics[width=\textwidth]{figures/fig2_temporal.png}
    \caption{\textbf{Figure 2 --- Temporal Patterns.} Left: Box plot of monthly temperatures showing seasonal variation. Right: Weekly mean temperature time series.}
    \label{fig:temporal}
\end{figure}

\subsection{Seasonal Cycle and Time Series}

As illustrated in the left panel, June--August exhibits the highest temperatures while November--January corresponds to the coldest months. Monthly temperature fluctuations range between 5 and 10 degrees Celsius. The median temperature difference between winter and summer is approximately $20\,^\circ\text{C}$, and the temperature gap between the coldest and hottest periods can exceed $30\,^\circ\text{C}$.

The right panel reveals a distinct seasonal cycle corroborating the monthly data. Temperature trends over the two-year period are relatively consistent, though minor deviations are observed---for instance, a notable temperature drop occurred around February 2015.

\subsection{Key Temporal Findings}


\newpage

% ══════════════════════════════════════════════════════════════
% Section 4: NMME Ensemble Models vs Target
% ══════════════════════════════════════════════════════════════
\section{NMME Ensemble Models vs Target}

\subsection{Individual Model Correlation}

\begin{figure}[H]
    \centering
    \includegraphics[width=\textwidth]{figures/fig3_nmme.png}
    \caption{\textbf{Figure 3 --- NMME Ensemble Models vs Target.} Left: Pearson correlation of individual NMME models with the target variable. Right: Scatter plot of NMME ensemble mean vs observed temperature, colored by season.}
    \label{fig:nmme}
\end{figure}

The left panel displays the correlation coefficients between each NMME model's forecasted 14-day 2\,m temperature and the observed values. The vast majority of models achieve correlations around $r \approx 0.90$, indicating strong predictive skill. More than 20 models exceed $r = 0.9$, including all CFSv2, GFDL, NASA, CanCM3/4, and most CCSM4 variants. However, the \textbf{CCSM3} models stand out with noticeably lower correlations, falling slightly below the 0.9 threshold, suggesting it is the least skillful ensemble member.

\subsection{Ensemble Mean Performance}

The right panel shows the ensemble mean plotted against observed temperatures, with points colored by season. The scatter points cluster tightly around the $y = x$ line, demonstrating excellent agreement with no systematic bias. The overall Pearson correlation is $r = 0.950$, confirming an extremely strong linear relationship. The four seasons are well mixed along the diagonal, indicating consistent performance across all seasons.

\newpage

% ══════════════════════════════════════════════════════════════
% Section 5: Feature Importance & Optimal Feature Count
% ══════════════════════════════════════════════════════════════
\section{Feature Importance \& Optimal Feature Count}

\subsection{Ensemble Feature Importance Ranking}

\begin{figure}[H]
    \centering
    \includegraphics[width=\textwidth]{figures/fig4_feature_importance.png}
    \caption{\textbf{Figure 4 --- Feature Importance \& Optimal Feature Count.} Left: Top 30 features ranked by ensemble importance score (Lasso + ElasticNet + Random Forest). Right: Learning curve of CV RMSE vs number of features.}
    \label{fig:feature_importance}
\end{figure}

The left panel reports the top 30 features based on the ensemble importance score, computed as the normalized average of Lasso, ElasticNet, and Random Forest importance. The most important feature is \texttt{contest-slp-14d\_\_slp}, followed by NMME temperature forecasts and other contest meteorological predictors. The top-10 features mainly come from \textbf{contest reanalysis variables} and \textbf{NMME climate forecasts}, suggesting that both short-term atmospheric conditions and subseasonal climate signals contribute substantially to 14-day temperature prediction. In addition, several SST and sea-ice features also appear in the top ranks---physically reasonable since oceanic boundary conditions can influence subseasonal temperature variability.

\subsection{Optimal Feature Count --- Learning Curve}

The right panel shows the learning curve of CV RMSE as features are added in order of importance. A clear elbow appears at around \textbf{242 features}, after which the RMSE shows minimal improvement. This indicates that the model has captured most useful information by this point, and additional low-importance features mainly introduce noise rather than improving predictive skill.